\documentclass{rvcepaper}
\begin{document}

% ---------- IM373TA Total Quality Management
\begin{iempaper}
\paperinfo{shownote=0,date={06/11/2025},exam={CIE -- I},maxmarks={10 + 50},sem={VII},prog={UG},duration={30 + 90 Min},
 title={Total Quality Management},code={IM373TA}}
\iemhead
{\small Note:}\par\vspace{-1pt}
\hspace*{0.5cm}{\footnotesize 1. \hspace{0.1cm}Answer all the Questions.}\par\vspace{3pt}
\begin{cietable}{M}{BT}{CO}
\qpart{Part -- A}
\q{1.}{From the Deming's management philosophy, identify the point which is relevant to managing supplier.}{1}{1}{CO1}
\q{2.}{\blank[2cm] is the title of the book authored by Philip B Crosby in 1979.}{1}{1}{CO1}
\q{3.}{Define fitness to use.}{2}{1}{CO1}
\q{4.}{Name the future fitness of quality which is futuristic.}{1}{2}{CO1}
\q{5.}{An employee is member of quality circle involved in quality problem solving. He has to analyse the relationship between two variables. Suggest a tool / tools for the purpose.}{2}{3}{CO2}
\q{6.}{List any two quality principles outlined by Juran.}{2}{2}{CO1}
\q{7.}{What is ISO 56000?}{1}{3}{CO1}
\qpart{Part -- B}
\q{1 a)}{Illustrate with a sketch the Juran's Trilogy for quality improvement.}{5}{2}{CO1}
\q{1 b)}{Elucidate Ishikawa's contribution to quality management.}{5}{2}{CO1}
\q{2}{Develop the sketch for evolution of methodology. List the practices and tools of quality used to achieve the fitnesses during the various stages of evolution.}{10}{3}{CO2}
\q{3}{Provide an outline of the Malcolm Baldrige National Quality Award (MBNQA)}{10}{2}{CO1}
\q{4}{Write a note on customer satisfaction. Give examples on how organisations achieve customer satisfaction.}{10}{3}{CO1}
\q{5}{Poor performance of students in examinations is the theme given to you. Explain briefly the application of 7QC steps for the above situation. Also list the 7QC tools that may be used.}{10}{4}{CO2}
\end{cietable}
\end{iempaper}

% ---------- IM374TFA Predictive Analytics
\begin{iempaper}
\paperinfo{shownote=0,date={07-11-2025},exam={CIE -- I},maxmarks={10 + 50},sem={7\textsuperscript{th}},prog={UG},duration={30 + 90 Min},
 title={PREDICTIVE ANALYTICS \ (Elective)},code={IM374TFA}}
\iemhead
{\small Note:}\par\vspace{-1pt}
\hspace*{0.5cm}{\footnotesize 1. \hspace{0.1cm}Answer all the Questions.}\par\vspace{3pt}
\begin{cietable}{M}{BT}{CO}
\qpart{Part -- A}
\q{1.}{Which type of data is most commonly mined---structured or unstructured?}{01}{L$_1$}{CO1}
\q{2.}{Name one common application of data mining in business.}{01}{L$_1$}{CO1}
\q{3.}{Name one method of data reduction.}{01}{L$_2$}{CO1}
\q{4.}{What is normalization in data preprocessing?}{01}{L$_2$}{CO1}
\q{5.}{What does EDA stand for?}{01}{L$_1$}{CO2}
\q{6.}{Which plot is best for visualizing data distribution?}{01}{L$_2$}{CO2}
\q{7.}{What type of regression is used for categorical outcomes?}{01}{L$_2$}{CO1}
\q{8.}{Which classification algorithm assumes feature independence?}{01}{L$_2$}{CO1}
\q{9.}{What is overfitting in model training?}{01}{L$_1$}{CO1}
\q{10.}{Which process ensures the model performs well on unseen data?}{01}{L$_2$}{CO1}
\qpart{Part -- B}
\q{1.}{What is Data Mining? What Is Predictive Analytics? Bring out the connection between data mining and predictive analytics with the help of examples. List any three areas of application for predictive analytics.}{10}{L$_3$}{CO1}
\q{2.}{List the five typical stages of a data mining project in their correct sequence. Discuss the relevance of CRISP-DM framework for carrying out data mining projects.}{10}{L$_3$}{CO1}
\q{3.}{Explain the primary goal of the \textbf{Data Preprocessing} stage in the data mining process and briefly describe the three main sub-tasks typically performed within it.}{10}{L$_3$}{CO2}
\q{4.}{Classify whether to `Play' based on Weather and Temperature.\newline Dataset:
\qtabc{|c|c|c|}{\hline \textbf{Weather} & \textbf{Temperature} & \textbf{Play}\\\hline Sunny & Hot & No\\\hline Sunny & Mild & No\\\hline Overcast & Hot & Yes\\\hline Rainy & Mild & Yes\\\hline Rainy & Cool & Yes\\\hline}
Apply any one of the classifier algorithms to come out with the final decision.}{10}{L$_4$}{CO2}
\q{5a.}{List the various types of classifiers commonly used in predictive analytics highlighting their main feature in one or two sentences.}{06}{L$_3$}{CO1}
\q{b.}{You are developing a system to classify job applicants in a management scenario as ``high potential'' or ``standard.'' Discuss the ethical considerations of using a highly accurate, but complex and non-interpretable model (like a Deep Neural Network) versus a slightly less accurate but easily interpretable model (like Logistic Regression or a Decision Tree) in this human resources context. Justify your choice of model type.}{04}{L$_4$}{CO2}
\end{cietable}
\stars{*\textasciitilde*\textasciitilde*}
\end{iempaper}

% ---------- IM374TFB Artificial Intelligence in Industrial and Systems Engineering
\begin{iempaper}
\paperinfo{shownote=0,date={7\textsuperscript{th} Nov 2025},exam={CIE-1},maxmarks={10 + 50},sem={VII},prog={UG},duration={30 + 90 Min},
 title={Artificial Intelligence in Industrial and Systems Engineering},code={IM374TFB}}
\iemhead
{\small Note:}\par\vspace{-1pt}
\hspace*{0.5cm}{\footnotesize 1. \hspace{0.1cm}Answer all the Questions.}\par\vspace{3pt}
\begin{cietable}{M}{BT}{CO}
\qpart{Part -- A}
\q{1}{Mention two key historical milestones in the evolution of AI.}{2}{3}{01}
\q{2}{What is the difference between Artificial Intelligence and Machine Learning?}{2}{2}{01}
\q{3}{How does data play a role in the functioning of a learning agent?}{2}{2}{02}
\q{4}{Mention two AI techniques commonly used in supply chain optimization.}{2}{3}{02}
\q{5}{Define demand forecasting in the context of supply chain management.}{2}{2}{02}
\qpart{Part -- B}
\q{1}{``The definition of Artificial Intelligence has evolved from mimicking human thought to achieving rational behavior.''\newline Explain this statement with reference to the four categories of AI definitions.}{10}{3}{01}
\q{2}{Explain how the Learning Element, Performance Element, Critic, and Problem Generator work together to enable learning and improvement in intelligent agents.}{10}{3}{02}
\q{3}{Explain how AI-enabled Industrial Engineering differs from traditional methods in terms of data usage, decision-making, and automation. Provide suitable industrial examples to support your answer.}{10}{3}{03}
\q{4}{``AI enables intelligent, self-optimizing supply chains.'' Examine this statement by describing how AI technologies such as machine learning, IoT, and robotics contribute to automation and adaptability in supply chain operations.}{10}{4}{03}
\q{5}{Discuss the key enablers of IE 4.0 and evaluate how Artificial Intelligence enhances industrial performance through smart data analysis, autonomous systems, and system-wide optimization}{10}{3}{03}
\end{cietable}
\end{iempaper}

% ---------- IM374TFD Sustainability Engineering & Circular Economy
\begin{iempaper}
\paperinfo{shownote=0,date={10\textsuperscript{th} Nov 2025},exam={CIE -- I},maxmarks={10 + 50},sem={VII},prog={UG},duration={30 + 90 Min},
 title={SUSTAINABILITY ENGINEERING \& CIRCULAR ECONOMY},code={IM374TFD}}
\iemhead
{\small Note: Answer all the questions}\par\vspace{3pt}
\begin{cietable}{M}{BT}{CO}
\qpart{Part -- A}
\q{1.}{What is the goal of SDG 13?}{02}{2}{1}
\q{2.}{What is meant by Sustainable Development?}{02}{2}{1}
\q{3.}{Mention the 3 Ps considered in Triple Bottom Line.}{02}{1}{1}
\q{4.}{What is Life Cycle Inventory in the context of Life Cycle Assessment?}{02}{2}{2}
\q{5.}{Mention the goal of Life Cycle Assessment.}{02}{1}{2}
\qpart{Part -- B}
\q{1.}{Explain the objectives of the Paris Agreement and critically analyze the role of engineering in achieving its climate goals.}{10}{4}{CO1}
\q{2.}{Discuss the key attributes expected of a graduate engineer that are essential for advancing sustainable engineering practices. Illustrate with relevant examples.}{10}{3}{CO1}
\q{3.}{Explain the concept of Life Cycle Thinking and its relevance to Circular Economy. Illustrate how engineers can apply these principles to promote sustainable industrial practices.}{10}{3}{CO1}
\q{4.}{Enumerate the steps involved in Life Cycle Impact Assessment (LCIA) and explain how each step contributes to evaluating the environmental performance of a product or process.}{10}{3}{CO2}
\q{5 a.}{Explain the iterative nature of Life Cycle Assessment (LCA) and its significance in improving the accuracy and reliability of sustainability evaluations.}{05}{2}{CO2}
\q{b.}{Discuss the relevance of Life Cycle Assessment (LCA) in engineering decision-making. Provide one example to illustrate its application.}{05}{3}{CO2}
\end{cietable}
\end{iempaper}

% ---------- IM375TGM Project Management (Institutional Elective)
\begin{iempaperB}
\paperinfo{ayear={2025-2026 (ODD Sem)},date={07.11.2025},code={IM375TGM},sem={VII Semester},exam={CIE -- I},maxmarks={10 + 50},duration={120 Mins},
 title={PROJECT MANAGEMENT (Institutional Elective)},shownote=1,hdr=1,usn=1,rule=1}
\iemheadB
\ptitle{PART -- A}
\begin{longtable}{|C{1.05cm}|p{12.2cm}|C{0.95cm}|C{0.95cm}|C{0.95cm}|}
\hline \textbf{Sl.\newline No.} & \multicolumn{1}{c|}{\textbf{Questions}} & \textbf{M} & \textbf{BT} & \textbf{CO}\\\hline
\q{1}{What is the relationship between project management and organizational strategy?}{02}{L2}{CO1}
\q{2}{List two techniques used for requirement collection.}{02}{L2}{CO1}
\q{3}{What is the key difference between projects and operations?}{02}{L1}{CO2}
\q{4}{By taking an example of real time project of your choice, identify any four success criteria of that project.}{02}{L1}{CO1}
\q{5}{State the purpose of feasibility analysis.}{02}{L2}{CO2}
\end{longtable}
\ptitle{PART -- B}
\begin{longtable}{|C{1.05cm}|p{12.2cm}|C{0.95cm}|C{0.95cm}|C{0.95cm}|}
\hline \textbf{Sl.\newline No.} & \multicolumn{1}{c|}{\textbf{Questions}} & \textbf{M} & \textbf{BT} & \textbf{CO}\\\hline
\q{1.}{As a newly appointed project manager in an engineering company, you are asked to explain how portfolio, program, and project management are interrelated and how they contribute to achieving organizational objectives.}{10}{L1}{CO1}
\q{2.}{As part of a knowledge-sharing workshop, you are explaining the PMBOK Guide to your cross-functional project team.
\begin{bul}\item Highlight the essential structure and components of the guide.\item Analyse how its application improves coordination among stakeholders and ensures timely project delivery.\end{bul}}{10}{L2}{CO2}
\q{3.}{An entrepreneur plans to invest in a renewable energy project. Discuss the steps involved in generation and preliminary screening of project ideas to ensure investment feasibility.}{10}{L3}{CO3}
\q{4.}{You are the project manager for developing an online hostel management system. Explain how you will collect and document the project requirements to ensure that all stakeholder needs are properly captured and understood.}{10}{L4}{CO4}
\q{5.}{Create a Work Breakdown Structure for developing a Chatbot to assist final-year students with queries related to placement schedules, company profiles, eligibility criteria, and interview guidelines at RV College of Engineering.}{10}{L2}{CO2}
\end{longtable}
\btfoot{|l|l|l|*{5}{C{0.8cm}|}*{6}{C{0.8cm}|}}{\hline
\multirow{2}{*}{Marks Distribution} & \multicolumn{2}{c|}{Particulars} & CO1 & CO2 & CO3 & CO4 & CO5 & L1 & L2 & L3 & L4 & L5 & L6\\\cline{2-14}
 & Max Marks & Quiz & 06 & 04 & -- & -- & -- & 04 & 06 & -- & -- & -- & --\\\cline{3-14}
 & & Test & 10 & 20 & 10 & 10 & -- & 10 & 20 & 10 & 10 & -- & --\\\hline}
\stars{*****}
\end{iempaperB}

% ---------- HS375TGT Principles & Practices of Cyber Law
\begin{iempaper}
\paperinfo{shownote=0,date={07/11/2025},exam={CIE - I},maxmarks={10 + 50},sem={VII},prog={UG},duration={30 + 90 Min},
 title={PRINCIPLES \& PRACTICES OF CYBER LAW},code={HS375TGT}}
\iemhead
{\small Note:}\par\vspace{-1pt}
\hspace*{0.5cm}{\footnotesize 1. \hspace{0.1cm}Answer all the Questions.}\par\vspace{3pt}
\begin{cietable}{M}{BT}{CO}
\qpart{Part -- A}
\q{1.}{Define Phishing}{2}{2}{CO1}
\q{2.}{What is the main purpose of the Information Technology Act, 2000 (IT Act)?}{2}{2}{CO3}
\q{3.}{Rajeev, a computer hacker, uses his skills to gain unauthorized access to a critical government website. Once inside, he uploads a malicious code that disrupts the website's operations, causing temporary paralysis in public services like online tax filings and public safety alerts. The government authorities discover that the attack was not only an act of hacking but was also intended to create widespread fear and panic among citizens. Based on the scenario above, explain whether the act qualifies as cyber terrorism under the Information Technology Act, 2000 (IT Act) and the punishment it may attract.}{2}{4}{CO2}
\q{4.}{Define Cyber Law.}{2}{2}{CO1}
\q{5.}{Mention any two preventive measures against cyber threats.}{2}{2}{CO1}
\qpart{Part -- B}
\q{1}{Amit is a software developer working in a reputed IT company in India. One day, he receives an email from an unknown source asking him to click on a link to update his work credentials. The link takes him to a page that looks like his company's login portal. He enters his username and password, but soon realizes that the website was fake and the email was a phishing attempt. By the time Amit tries to secure his account, someone has already used his credentials to access confidential company data and cause a breach in security. Based on the scenario above, explain the type of cybercrime committed and discuss the relevant provisions of the Information Technology Act, 2000 (IT Act) that apply in this case.}{10}{4}{CO1}
\q{2}{How do the challenges faced by law enforcement agencies in handling cybercrimes differ from those encountered in dealing with conventional crimes explain with one case example.}{10}{2}{CO2}
\q{3}{Discuss the various types of cybercrimes prevalent in India and explain the measures individuals can take to safeguard themselves against such crimes.}{10}{3}{CO1}
\q{4}{Explain the concept of Cyber Jurisdiction. Discuss its nature, scope, and challenges. Support your answer with reference to Indian and international perspectives with relevant judicial interpretations.}{10}{3}{CO2}
\q{5}{Riya, a college student, discovers that her photographs have been digitally altered (morphed) and circulated on social media platforms without her consent. The images were modified in a sexually explicit manner and shared widely through fake profiles. Distressed, she approaches the Cyber Crime Cell to file a complaint under the Information Technology Act, 2000.\newline
Answer the following questions:\newline
a) What is morphing, and how does it amount to a cybercrime under Indian law?\newline
b) Which provisions of the Information Technology Act, 2000 and the Indian Penal Code (IPC) are applicable to this case?\newline
c) Explain the legal remedies available to Riya.\newline
d) Mention any relevant case laws decided under similar circumstances.\newline
e) Discuss the role of intermediaries (such as social media platforms) in preventing the circulation of such content.}{10}{3}{CO2}
\end{cietable}
\end{iempaper}

\end{document}
